% Corrigé intégré automatiquement pour la banque Exercices.
% Source : A_Integrer/DDS_04/083_SchemaBlocs_FT/083_SchemaBlocs_FT.tex
\begin{corrigebox}[Corrigé]
\CorrigeQuestion{1}
On réduit le schéma-blocs de l'intérieur vers l'extérieur : associations en série par produit, associations en parallèle par somme et boucle de retour par $H_{BF}=\dfrac{G}{1+GH}$ (retour négatif). Après simplification, on ordonne le numérateur et le dénominateur selon les puissances de $p$, puis on identifie le gain statique, la classe et les paramètres canoniques.
\CorrigeQuestion{2}
On définit les variables, le pas de calcul et les conditions initiales, puis on traduit l'équation par une relation de récurrence. L'algorithme est vérifié sur les premiers pas et la stabilité numérique est contrôlée en comparant le pas aux constantes de temps du système.
\CorrigeQuestion{3}
On réduit le schéma-blocs de l'intérieur vers l'extérieur : associations en série par produit, associations en parallèle par somme et boucle de retour par $H_{BF}=\dfrac{G}{1+GH}$ (retour négatif). Après simplification, on ordonne le numérateur et le dénominateur selon les puissances de $p$, puis on identifie le gain statique, la classe et les paramètres canoniques.
\CorrigeQuestion{4}
La résolution consiste à identifier les données et l'inconnue, choisir la loi du cours adaptée, écrire l'équation symbolique avant toute application numérique, puis vérifier l'homogénéité, le signe et l'ordre de grandeur du résultat. Les hypothèses utilisées doivent être explicitement contrôlées à la fin.
\CorrigeQuestion{5}
La résolution consiste à identifier les données et l'inconnue, choisir la loi du cours adaptée, écrire l'équation symbolique avant toute application numérique, puis vérifier l'homogénéité, le signe et l'ordre de grandeur du résultat. Les hypothèses utilisées doivent être explicitement contrôlées à la fin.
\CorrigeQuestion{6}
On définit les variables, le pas de calcul et les conditions initiales, puis on traduit l'équation par une relation de récurrence. L'algorithme est vérifié sur les premiers pas et la stabilité numérique est contrôlée en comparant le pas aux constantes de temps du système.
\CorrigeQuestion{7}
La résolution consiste à identifier les données et l'inconnue, choisir la loi du cours adaptée, écrire l'équation symbolique avant toute application numérique, puis vérifier l'homogénéité, le signe et l'ordre de grandeur du résultat. Les hypothèses utilisées doivent être explicitement contrôlées à la fin.
\CorrigeQuestion{8}
On applique le théorème de la valeur finale : $y(\infty)=\lim_{p\to0}pY(p)$, après avoir vérifié la stabilité de $pY(p)$. L'écart permanent se déduit de $E(p)=\dfrac{R(p)}{1+L(p)}$ en tenant compte de la classe de la boucle ouverte et de la nature de l'entrée.
\end{corrigebox}
